% =====================================================================
% 4. Architectures in PyTorch (MLLB, Maria Martinez)
% Own blocks in depth: src/models/rnn.py (T11), src/models/attention.py and
% src/models/transformer.py (T12). Dense and CNN (Oriana) summarised.
% =====================================================================
\section{Architectures in PyTorch}\label{sec:mllb-arch}

This section explains how my models are implemented, with the PyTorch details that matter for them; the
theory and the analysis of the results are in my PRDL report. All models map (batch, 12, 1000) to five
logits, and my models end with the head of the CNN baseline: dropout 0.5 and one linear layer.

\subsection{Dense network and CNN (Oriana's blocks)}

Oriana's dense network flattens the record (\texttt{nn.Flatten}) and applies fully connected layers with
batch norm, ReLU and dropout. Her CNN stacks residual blocks of \texttt{Conv1d}, batch norm and ReLU with a
$1\times1$ convolution on the shortcut when the number of channels changes, and ends with global average
pooling. I reused the first three of her convolutional blocks, unchanged, as the front end of my best
recurrent model.

\subsection{Recurrent networks}

\paragraph{From a signal to a sequence.} \texttt{nn.LSTM(batch\_first=True)} expects (batch, time,
features), the opposite axis order of \texttt{Conv1d} (batch, channels, time). Listing~\ref{lst:rnnseq}
builds the sequence. For frames, \texttt{reshape(b, c, t, frame)} cuts each lead into $t$ frames,
\texttt{permute(0, 2, 1, 3)} moves time before the leads, and the last \texttt{reshape} flattens the 12
leads $\times$ 10 samples of each frame into one vector of 120 numbers; a test checks this element by
element. For the convolutional front end, the output of the CNN blocks (batch, 128, 125) is simply
permuted to (batch, 125, 128). No padding is needed, because every record has exactly 1000 samples.

\lstinputlisting[firstnumber=152,caption={From the signal to the sequence read by the RNN, and the final state (\texttt{src/models/rnn.py}).},label={lst:rnnseq}]{common/code/rnn_sequence_encode.py}

\paragraph{What \texttt{nn.LSTM} returns.} The second part of Listing~\ref{lst:rnnseq} handles the outputs:
\begin{itemize}[leftmargin=1.2em,itemsep=1pt]
  \item \texttt{out} has the hidden states of the top layer at every step, (batch, $T$, hidden) or
        (batch, $T$, 2 $\times$ hidden) if bidirectional.
  \item The LSTM also returns \texttt{(h\_n, c\_n)}, the final hidden and cell states of every layer and
        direction; the GRU returns only \texttt{h\_n}, because it has no cell state. The code checks for a
        tuple so that one class handles both cells.
  \item For a unidirectional network, \texttt{h\_n[-1]} (top layer) equals \texttt{out[:, -1, :]}.
  \item For a bidirectional network the final summary is the concatenation of \texttt{h\_n[-2]}, the
        forward direction, which has read the record from left to right, and \texttt{h\_n[-1]}, the
        backward direction, which has read it from right to left. The backward final state is at
        \texttt{out[:, 0, hidden:]}, not at \texttt{out[:, -1, :]}: taking the last row of \texttt{out}
        would pair the forward summary with a backward state that has read only one step. A test checks
        both facts.
\end{itemize}
Two more details. The \texttt{dropout} argument of \texttt{nn.LSTM} only acts between stacked layers, so it is
left at 0 and dropout is applied to the summary vector before the classifier. And a test checks causality:
in a unidirectional network, changing the second half of the record leaves the states of the first half
unchanged.

\paragraph{Attention pooling.} Listing~\ref{lst:attnpool} computes the scores, the weights and the context
vector. \texttt{nn.Linear} acts on the last axis, with the same weights at every time step, so it is already
``time distributed''. The softmax is taken over time (\texttt{dim=1}), and \texttt{torch.bmm} multiplies the
weights (batch, 1, $T$) by the states (batch, $T$, $d$) to give the context vector (batch, 1, $d$). The weights
are returned too, for the figures.

\lstinputlisting[firstnumber=66,caption={Attention scores and context vector over the RNN states (\texttt{src/models/rnn.py}).},label={lst:attnpool}]{common/code/attention_pooling.py}

\paragraph{Gradients with respect to the input.} The sensitivity analysis needs
$\partial\,\mathrm{logit}/\partial x$: the input tensor gets \texttt{requires\_grad\_(True)} and
\texttt{torch.autograd.grad} computes the gradient of the summed logit. In eval mode the records of a batch do
not interact, so the gradient of the sum gives each record's gradient. This is computed on the CPU, because
cuDNN only runs the backward pass of an RNN in training mode, where dropout would change the result.

\subsection{Self-attention and the Transformer, written by hand}

I wrote multi-head self-attention myself instead of using \file{nn.MultiheadAttention} or
\file{nn.TransformerEncoderLayer}, so that every step can be explained and the attention weights are
available; a test checks that my module gives the same output as \texttt{nn.MultiheadAttention} with the
same weights, with and without a causal mask.

\paragraph{Patches as tokens.} Listing~\ref{lst:patches} cuts the record into patches with the same
reshape and permute as the RNN frames, so both models see the same 100 steps. One linear layer embeds each
patch; it is exactly a \texttt{Conv1d} with kernel and stride equal to the patch length (checked by a test).
An optional CLS token is an \texttt{nn.Parameter} of shape (1, 1, $d$), expanded to the batch and
concatenated in front of the patches. The sinusoidal positional encoding is registered as a buffer with
\texttt{persistent=False}: it moves to the GPU with the model, the optimiser does not change it, and it is
not saved in the checkpoint, because it is rebuilt when the model is created.

\lstinputlisting[firstnumber=96,caption={Patch tokens, CLS token and positional encoding (\texttt{src/models/transformer.py}).},label={lst:patches}]{common/code/patch_tokens.py}

\paragraph{Scaled dot-product attention.} In Listing~\ref{lst:attn}, \texttt{q @ k.transpose(-2, -1)} computes
all query--key dot products at once, shape (\ldots, $T$, $T$); they are divided by $\sqrt{d_k}$, the dimension
of one head, not of the model. Masked scores are set to $-\infty$ with \texttt{masked\_fill} before the
softmax, so their weights are exactly 0; the softmax is over the keys (last axis), so each query's weights
sum to 1. My classifier does not use the mask, but Oriana's generator reuses this same function with the
causal mask of \texttt{causal\_mask}, a lower-triangular boolean matrix.

\lstinputlisting[firstnumber=29,caption={Causal mask and scaled dot-product attention (\texttt{src/models/attention.py}).},label={lst:attn}]{common/code/causal_attention.py}

\paragraph{Splitting into heads.} Listing~\ref{lst:mhsa} projects the input into queries, keys and values
with three bias-free linear layers and splits $d_\text{model}$ into heads: \texttt{view(b, t, heads,
head\_dim)} cuts the last axis and \texttt{permute(0, 2, 1, 3)} moves the heads before time, giving (batch,
heads, $T$, head\_dim), so the attention runs on every head in parallel. Joining the heads back needs
\texttt{reshape}, not \texttt{view}: after \texttt{permute} the tensor is no longer contiguous in memory and
\texttt{view} would fail. $W_O$ then mixes the concatenated heads.

\lstinputlisting[firstnumber=93,caption={Multi-head self-attention: heads, attention and $W_O$ (\texttt{src/models/attention.py}).},label={lst:mhsa}]{common/code/mhsa_split_forward.py}

\paragraph{Pre-LN and Post-LN in one function.} Listing~\ref{lst:sublayer} applies one sub-layer with its
skip connection and LayerNorm. With Pre-LN the input is normalised before the sub-layer and the result is
added to the unnormalised residual stream; with Post-LN the LayerNorm is applied after the addition. The
flags \texttt{residual}, \texttt{ffn} and \texttt{out\_proj} remove components, which is how the five models of
the step-by-step study are built from the same class. Because Pre-LN leaves the residual stream
unnormalised after the last block, the encoder applies a final LayerNorm. The attention maps are returned by
every block and collected in eval mode on the CPU.

\lstinputlisting[firstnumber=173,caption={One sub-layer with skip connection and LayerNorm (\texttt{src/models/attention.py}).},label={lst:sublayer}]{common/code/block_sublayer.py}

\subsection{Results of my blocks in brief}

The final models are an LSTM with a convolutional front end (\texttt{conv},
\rnnFinalTestMacroAurocMean{}~$\pm$~\rnnFinalTestMacroAurocSd{} on test), an LSTM with attention over 1000 raw
steps (\texttt{raw\_attn}, \rnnAttnfinalTestMacroAurocMean{}~$\pm$~\rnnAttnfinalTestMacroAurocSd{}) and a
two-block Transformer without positional encoding (\texttt{no\_pos},
\transformerFinalTestMacroAurocMean{}~$\pm$~\transformerFinalTestMacroAurocSd{}), against
\cnnFinalTestMacroAurocMean{} for the CNN (Table~\ref{tab:comparison_final}). The implementation points that
mattered most for these results were the axis order of the sequence (any mistake there would silently mix
leads and time), reading all hidden states instead of the last one, and the shapes of the attention heads.
